% !TeX root = ../papers/nc-p3.tex
\section{種数三の母線写像の平滑化と孤立三点繊維}
\label{sec:genus-three-kernel-excess}
\begin{theorem}\label{thm:genus-three-kernel-excess}
被約四辺有限標準正則窓の整反対八次で余因子のスキーム次数が十八なら核像は曲面である。
一般階数一曲線が存在しても成立する。
\end{theorem}
\begin{proof}[種数三の底への降下]
曲線核 $K$ を仮定する。前定理と定理\ref{thm:sectional-cofactor-length}より、
一般平面の自己準同型曲線 $G$ は整Gorenstein、算術種数三、核束 $P$ は可逆で次数六、
$h^0(P)=4$。$G\to K$ は有限双有理である。実際入力の正規化母線束上で一般平面は
$K$ の正規化の一切断を与える。Castelnuovo上界から $3\le p_a(K)\le4$。
三なら $K=G$、四なら $K$ は二次・三次の完全交叉でGorensteinである
（定理\ref{thm:cofactor-curve-kernel-only-bezout}の六次種数四の議論）。
後者で $G$ は $K$ の一点における余長一の部分正規化であり、
各特異点でそのような中間環は $\mathfrak m^{-1}$ しかない。
従って同定理の有限候補と構成可能性の議論により、一つの固定曲線 $G$ に対して
一般入力平面との交点写像が降りる。二つの一般平面を使うと母線写像も $G$ に降りる。
$G$ 上の母線部分束 $U\subset V\otimes\mathcal O_G$ の商を $F$ とすれば
$\deg F=8$。射影直線束の入力射 $f$ は有限で、Euler多項式の比較により
$M_D=f_*\pi^*P$ である。実際右辺は $(t+1)(4t+6+1-3)$、核座標が一般同型な評価包含を与える。
定理\ref{thm:genus-two-kernel-excess}と同様に少なくとも十四個の孤立三点繊維を持つ。
また $F\to P$ は余核の一次関係を与えるから $\operatorname{Hom}(F,P)=0$。

$F$ は半安定である。次数三以下の階数一torsion-free商 $D'$ があれば
$\mathcal Hom(D',P)$ の次数は $6-\deg D'\ge3$ なのでRiemann--Rochで非零切断を持つ。
ここでGorenstein双対性から階数一torsion-free層の双対の次数は元の次数の負である。
よって $F\otimes\omega_G^{-1}$ は傾き零の半安定束。
これから $U$ への非零射があれば、像は次数非負で $V\otimes\mathcal O_G$ の部分層となる。
その階数と同じ本数の一般成分に射影すれば自明束に一般同型に入り、次数比較で同型である。
従って元の自明四束での像も定数部分空間で、全母線に共通入力点ができる。
これは核曲線の非退化性と $N$ の成分の基点自由性に反する。
従って $H^1(U^*\otimes F)=0$ である。
\end{proof}

\begin{proof}[局所平面性と滑らかな底上の次元]
$G$ は局所平面曲線である。正規化余長は各点で $\delta\le3$。
非平面点なら、一般の一次パラメータで局所環 $T$ を割ったArtin Gorenstein環は
埋込次元二以上、従って長さ $m\ge4$ である（極大イデアルの平方が零ならsocleが二次元以上）。
正規化環 $S$ でこのパラメータを $x$ とすると極大イデアルは $xS$ を生成し、
$T\subset k+xS$ より $\delta\ge m-1$。
$\delta\le3$ は $m=4$, $T=k+xS$ を強制するが、その導手は $xS$ であり
Gorenstein曲線の導手公式 $\operatorname{length}(S/(T:S))=2\delta$ に反する。
この公式は有限双対で $\operatorname{length}(T/(T:S))=\operatorname{length}(S/T)$ から従う。
ゆえに命題\ref{prop:cofactor-kernel-smoothing}の局所平面曲線の平滑化と
Grassmann射の持上げを適用できる。

滑らかな種数三曲線上で、この射を三点繊維が有限な射に変形できることを示す。
固定した底でのGrassmann射の空間は、$F$ が半安定で共通入力点がない開集合では
前段と同じ消滅により滑らかで、次元は
$\chi(U^*\otimes F)=32-4\cdot2=24$。
束と順序付き四生成切断を指定し、束の自己同型で割って数える。
$h^1(F)>0$ なら半安定性から $F\twoheadrightarrow\omega$ があり、核は次数四の線束 $C$。
$C\ne\omega$ なら $h^1(F)=1$、$C$ は三次元を動き、非分裂拡大の射影空間は一次元。
この族は四切断とスカラー自己同型を数えても $3+1+20-1=23$ 次元以下である。
分裂の場合もこれ以下。$C=\omega$ の非分裂拡大は射影平面で $h^1(F)=1$、
分裂束は $h^1=2$ だが自己同型次元四なので高々 $24-4=20$ 次元である。
よって一般には $h^1(F)=0$, $h^0(F)=4$。
この開部分で厳密半安定の束は次数四線束二つの拡大であり、
異なる二束では二つのPicardの六次元と一次元の射影拡大、
同じ束では三次元と二次元の射影拡大だけである。分裂の場合も含め、
四切断の $16$ 次元とスカラー自己同型の分を加えても $24$ に達しない。
従って各該当成分で一般の束は安定となる。

安定束で三点以上の零を持つ切断の飽和線部分束 $C$ は次数ちょうど三である。
商は次数五の線束 $D'$。線部分束を印した拡大の族の次元は二つのPicardの $3+3$ と
$\mathbb P\operatorname{Ext}^1(D',C)$ の三次元、さらに四切断と自己同型の分 $16-1$ を足し、
高々二十四である。
従って各二十四次元成分の一般点ではこのような部分束の数は有限である
（支配しない入射成分は一般点を避ける）。また $h^0(C)\ge2$ なら
Riemann--Rochより $C=\omega(-p)$ で、最初の三次元は一次元に下がる。
従って一般の部分束は各々高々一つの射影切断を与え、三点繊維は有限である。
以上の次元評価は拡大を線束のPicardと相対拡大空間でパラメータ化するので有限型の族の評価である。
半安定性と共通入力点のない条件は平滑化でも開であり、
相対Homスキームにおける閉包と曲線の選択は定理\ref{thm:genus-two-kernel-excess}と同じである。
\end{proof}

\begin{proof}[三点繊維の長さ]
元の十四個以上の孤立点は局所平面曲線上の長さ三部分スキームなので、
同定理の相対Hilbertスキームの次元評価で各々が一般繊維に寄与する。
しかし滑らかで有限な一般繊維での長さは系\ref{cor:kernel-curve-triple-ruling-genus}から
$20-4\cdot3=8$ である。この矛盾で曲線核は排除される。
\end{proof}
\path{computations/verify_genus_three_kernel_excess.py} は各拡大族とGrassmann射の次元、
三点繊維の長さを検算する。
